% ===========================================================================
% Section 1 -- Introduction: Can Machines Invent?
% File: sections/01_introduction.tex
% Reader question: Can an AI create knowledge no human wrote down?
% One-sentence answer: Skeptics say models trained on history can only
%   recombine it. AlphaGo Zero showed that a learner with a referee can go
%   beyond history, and we test that for invention.
% Core points advanced: 4 (skeptics right about imitation, wrong about
%   learning), 3 (search = 1), 1 (innovation is a game).
% Exhibits: Figure 1 (front matter, label fig:emptysquare) and Box 1.
% Budget: 1.5 pages. Editor notes: sections/01_introduction.notes.md
% ===========================================================================

\section{Can Machines Invent?}\label{sec:intro}

Can a machine create knowledge that no human has written down? The strongest
skeptics say no: a model trained on the record of human knowledge can only
recombine it. Go suggests another answer. AlphaGo Zero, taught by a referee
instead of a record, went beyond the best human play without seeing a human
game \citep{silver2017mastering}. We test whether invention
can do the same.

\paragraph{The skeptics' challenge.}
In March 2025, Yann LeCun argued that language models retrieve existing
solutions rather than find new ones: ``in their purest form, they are incapable
of inventing new things'' \citep{lecun2025bigtech}. The same month, Thomas Wolf
warned that current trends will give us ``a country of yes-men on servers''
\citep{wolf2025einstein}. Dwarkesh Patel had asked Dario Amodei in August 2023
what to make of models with ``basically the entire corpus of human knowledge
memorized'' that had not made ``a single new connection that has led to a
discovery'' \citep{patel2023question}. In January 2025, Demis Hassabis posed
the hardest version as a benchmark for general intelligence: today's systems
``can play a game of Go at a world champion level. But could a system invent
Go?'' \citep{hassabis2025bigtech}. The skeptics have found a real limit. A
model trained to imitate inherits the judgment of the record it imitates, so
its sense of what will work next is history's sense. That is a claim about the
training signal, not the network (\cref{sec:anatomy}).

\paragraph{Humans are bounded by history too.}
Research effort keeps rising while its productivity
falls: Moore's law now needs more than 18 times as many researchers as in the early 1970s \citep{bloom2020ideas}. Under a growing burden
of knowledge, inventors reach the frontier later and specialize more narrowly
\citep{jones2009burden}. Papers and patents have grown less likely to
break with the past \citep{park2023disruptive}, and fields open to outsiders
when their leading scientists die \citep{azoulay2019funeral}. The bottleneck is
not a shortage of minds. It is the gravity of what we already know.

\paragraph{Everyone is scaling search.}
The field's answer to both limits has been more tries. Andrej Karpathy's
autoresearch, released in March 2026, runs about 100 training experiments
overnight and keeps each change only if it helps
\citep{karpathy2026autoresearch}. The AI Scientist \citep{lu2024aiscientist}
and the AI co-scientist \citep{gottweis2025coscientist} run agentic research
loops, and AlphaEvolve evolves programs against automated evaluators
\citep{novikov2025alphaevolve}. Richard Socher's \emph{Eureka Machine} imagines
``millions of AI agents'' that ``critique each other's ideas''
\citep{socher2026bookbite,socher2026eureka}, and his company, Recursive, raised
\$650 million to automate the research cycle \citep{recursive2026techcrunch}.
These systems work, and their evaluators are exactly the referees our game
needs. But each pays for judgment at test time, one try at a time, and pays
again on the next problem. What a search learns stays mostly in its outputs,
not in the model that proposes the next idea.

\begin{keyidea}
Everyone else is buying judgment by the try. We learn it once.
\end{keyidea}

\paragraph{The Zero move.}
Go faced the same choice and took the other road. AlphaGo learned from human
expert games and searched heavily at every move \citep{silver2016mastering}.
AlphaGo Zero learned only from games against itself, and its network alone,
with no search, rated 3,055 Elo. That is close to the 3,144 of AlphaGo Fan, the
version that beat the European champion Fan Hui
\citep{silver2016mastering,silver2017mastering}. Search had been distilled into
intuition. Can invention have its own Zero moment? Make innovation a game,
learn judgment once by self-play, then invent in one move: \searchone{}
(\hyperref[box:searchone]{Box~1}). \Cref{fig:emptysquare} places this in the
one square that Go has filled and invention has not.

\begin{explainer}{Box 1: What ``search = 1'' means}
\phantomsection\label{box:searchone}%
\begin{itemize}[leftmargin=*,nosep]
  \item \textbf{One inference, one candidate.} For a new problem, the program
    writer emits one complete program ($B = 1$), with no retries, repairs,
    best-of-$N$, reranking or weight updates.
  \item \textbf{Everything is counted,} including reasoning tokens and problem
    encoding. One inference is one generation, not one forward pass.
  \item \textbf{Fitting is not search.} The program is fit to the training
    data like any learning algorithm; fitting sets parameters, never chooses
    programs. Training search is priced in \cref{sec:price}.
\end{itemize}
\textbf{Zero search, not zero knowledge:} the writer starts from pretraining and
human ideas.
\end{explainer}

\paragraph{What we show.}
We make innovation a game: the encoded problem is the board, a complete program
is the move, a verifier is the referee and the best known algorithm is the
opponent (\cref{sec:game}). \EZ{} plays it with a program writer that
proposes moves and a problem--program map that judges them, like AlphaGo's
policy and value networks, trained by self-play (\cref{sec:method}). The
contest is narrow by design. On \TBD{M} unseen problems spanning binary and
multiclass classification and regression, one inference \TBD{beats} the
algorithms that Claude, Codex and Gemini invented when asked for a generally
competitive machine-learning method,\footnote{Model versions, dates, the
verbatim prompt and every invented algorithm are published in the appendix.}
winning \TBD{K} of \TBD{M} (\cref{tab:headtohead}). It also \TBD{beats}
the foundation model KumoRFM \citep{kumorfm2025}, at an Elo of \TBD{E} (\cref{fig:leaderboard}). The best frontier model, run as an
agent with validation feedback, needs \TBD{R} tries to match it
(\cref{fig:searchone}). The problem--program map cuts wasted moves from
\TBD{X}\% to \TBD{Y}\% (\cref{tab:ablation}). Our contributions are five:
\begin{enumerate}[leftmargin=*,nosep,label=\arabic*.]
  \item \textbf{The anatomy of innovation:} four pillars (aim, variation,
    judgment and retention), and the observation that search length is the
    reciprocal of taste, judgment learned in advance (\cref{sec:anatomy}).
  \item \textbf{The innovation game:} rules that make invention playable on
    many boards (\cref{sec:game}).
  \item \textbf{The economics of algorithmic self-play,} from 8B to 1.5T
    parameters (\cref{sec:price}).
  \item \textbf{\EZ{}:} a program writer and a problem--program map that learn
    taste (\cref{sec:method}).
  \item \textbf{Evidence:} one inference against frontier-invented algorithms
    and KumoRFM on unseen problems, and the tries an agent needs to match it
    (\cref{sec:results}).
\end{enumerate}
\Cref{sec:road} traces each pillar's history in AI; \cref{sec:meaning} asks
what changes and where the claim stops.

To see why a referee changes everything, we first ask what innovation is made
of.
